\documentclass{beamer}
\setbeamerfont{headline}{size=\footnotesize}

\mode<presentation>{
\usetheme{Montpellier}
\setbeamercovered{transparent}
\usecolortheme{lsc}
}

\mode<handout>{
  \usepackage[bar]{beamerthemetree}
  \beamertemplatesolidbackgroundcolor{black!5}
}

\usepackage{amsmath,amssymb}
\usepackage[brazil]{varioref}
\usepackage[english,brazil]{babel}
\usepackage[utf8]{inputenc}
\usepackage{fontawesome}
\usepackage{tikz}
\usepackage{graphicx}
\usepackage{url}
\usepackage{colortbl}
\usepackage{caption}
\usepackage{transparent}
\usepackage{multirow}
\usepackage{listings} % Para código


% Configurações para o bloco de código
\lstset{
  basicstyle=\ttfamily\small,
  commentstyle=\color{gray},
  keywordstyle=\color{blue},
  stringstyle=\color{red},
  showstringspaces=false,
  breaklines=true,
  frame=single,
  framesep=5pt,
  backgroundcolor=\color{white},
  numbers=left,
  numberstyle=\tiny\color{gray},
  tabsize=2,
  columns=fullflexible,
  literate={á}{{\'a}}1 {é}{{\'e}}1 {í}{{\'i}}1 {ó}{{\'o}}1 {ú}{{\'u}}1
           {Á}{{\'A}}1 {É}{{\'E}}1 {Í}{{\'I}}1 {Ó}{{\'O}}1 {Ú}{{\'U}}1
           {ã}{{\~a}}1 {õ}{{\~o}}1 {ç}{{\c{c}}}1 {Ç}{{\c{C}}}1
           {à}{{\`a}}1 {è}{{\`e}}1 {ì}{{\`i}}1 {ò}{{\`o}}1 {ù}{{\`u}}1
}

% TikZ libraries for drawing
\usetikzlibrary{arrows.meta, positioning, shapes.geometric, calc, shadows, decorations.pathreplacing,matrix, arrows.meta, positioning, chains}

\beamertemplatetransparentcovereddynamic
\addtobeamertemplate{title page}{\centering{ \includegraphics[width=1.25cm]{figs/logo.png}}}{}
\title[Estrutura de Dados I]{Ordenação}
\author[Elton Sarmanho]{
Professor: Elton Sarmanho\inst{1}
\newline
E-mail: \href{mailto:eltonss@ufpa.br}{eltonss@ufpa.br}\newline
\href{https://github.com/eltonsarmanho}{\faGithub} \href{https://www.linkedin.com/in/elton-sarmanho-836553185/}{\faLinkedinSquare}
}
\institute[UFPA]{
   \inst{1}%
     Faculdade de Sistemas de Informação - UFPA/CUTINS
     }

\date{\today}
\pgfdeclareimage[height= 1.25 cm]{inf}{figs/logo.png}
\logo{{\transparent{0.35}\pgfuseimage{inf}}}

\expandafter\def\expandafter\insertshorttitle\expandafter{%
  \insertshorttitle\hfill%
  \insertframenumber\,/\,\inserttotalframenumber}

% TikZ Styles
\tikzstyle{arraynode}=[draw, minimum size=0.7cm, font=\bfseries]
\tikzstyle{highlight}=[fill=blue!30]
\tikzstyle{sorted}=[fill=green!30]
\tikzstyle{pivot}=[fill=orange!40]
\tikzstyle{pointer}=[->, >=Latex, thick, blue!80]
\tikzstyle{info_text} = [font=\small\itshape, text width=3cm]

\begin{document}
\setbeamertemplate{logo}{}
\begin{frame}
\titlepage
\end{frame}
\logo{{\transparent{0.35}\pgfuseimage{inf}}}

\begin{frame}
\frametitle{Roteiro}
\tableofcontents[sections={1-2}]
\end{frame}

\begin{frame}
\frametitle{Roteiro}
\tableofcontents[sections={3-4}]
\end{frame}

%%%%%%%%%%%%%%%%%%%%%%%%%%%%%%%%%%%%%%%%
\section{Conceitos Gerais}
\subsection{Objetivos}
\begin{frame}{Objetivos da Aula}
\begin{itemize}
 \item Esta aula apresenta os conceitos fundamentais de ordenação e os principais algoritmos descritos na literatura.
 \item Ao final, você deverá ser capaz de:
 \begin{itemize}
  \item Compreender o processo de comparação e movimentação de elementos em um algoritmo de ordenação.
  \item Analisar e comparar o desempenho dos diferentes algoritmos.
  \item Identificar aplicações reais onde a ordenação é fundamental.
 \end{itemize}
\end{itemize}
\end{frame}

\subsection{Definição}
\begin{frame}{Ordenação, Classificação e Organização}
\begin{itemize}
 \item Um \textbf{algoritmo de ordenação} reorganiza uma coleção de itens em uma ordem específica (crescente ou decrescente). A entrada é uma lista, e a saída é uma permutação dessa mesma lista, mas ordenada.
 \item A organização de dados é um pilar em ciência da computação:
 \begin{itemize}
  \item Algoritmos de busca (como a busca binária), fusão de listas e outros processos dependem de dados previamente ordenados.
  \item Aplicações em bancos de dados, geometria computacional e bioinformática necessitam de listas ordenadas para operar eficientemente.
 \end{itemize}
 \item Para humanos, dados ordenados são mais fáceis de ler, interpretar e analisar.
\end{itemize}
\end{frame}

\begin{frame}{Condições de Saída}
\begin{itemize}
 \item A saída de um algoritmo de ordenação deve satisfazer duas condições cruciais:
 \begin{enumerate}
  \item \textbf{Ordem:} A lista final deve estar em uma ordem específica (ex: não decrescente).
  \item \textbf{Permutação:} A lista final deve conter exatamente os mesmos elementos da lista original, apenas reorganizados.
 \end{enumerate}
\end{itemize}
\begin{center}
\begin{tikzpicture}[scale=0.8, transform shape]

    % Unsorted Array
    \foreach \val/\idx in {6/1, 5/2, 7/3, 1/4, 4/5, 3/6, 2/7} {
      \node[arraynode] (unsorted-\idx) at (\idx, 0) {\val};
    }
    % Rótulo para a entrada (acima do primeiro elemento)
    \node[info_text, above=0.5cm of unsorted-1] {Entrada (não ordenada):};

    % Arrow
    \draw[->, very thick, blue!60, decorate, decoration={snake,amplitude=.4mm,segment length=2mm,post length=1mm}] (4, -0.5) -- (4, -2);

    % Sorted Array
    \foreach \val/\idx in {1/1, 2/2, 3/3, 4/4, 5/5, 6/6, 7/7} {
      \node[arraynode, sorted] (sorted-\idx) at (\idx, -2.5) {\val};
    }
    % Rótulo para a saída (abaixo do primeiro elemento ordenado)
    \node[info_text, below=0cm of sorted-1] {Saída (ordenada):};

\end{tikzpicture}
\end{center}
\end{frame}

\begin{frame}{Motivação}
\begin{itemize}
 \item O estudo de algoritmos de ordenação é fundamental para a formação em computação:
\begin{itemize}
 \item Solidifica conceitos teóricos importantes:
 \begin{itemize}
  \item Notação Assintótica (Big O).
  \item Análise de pior caso, melhor caso e caso médio.
  \item O balanço (trade-off) entre tempo de execução e uso de memória.
 \end{itemize}
 \item Como a ordenação é um passo intermediário em muitos algoritmos complexos, é crucial conhecer os métodos mais eficientes para cada cenário.
\end{itemize}
\end{itemize}
\end{frame}

\subsection{Terminologia}
\begin{frame}{Terminologia Essencial}
\begin{itemize}
 \item \textbf{Entrada}: Uma sequência de $n$ itens, como $r[0], r[1], ..., r[n-1]$.
 \item \textbf{Registro}: Cada item individual na coleção a ser ordenada.
 \item \textbf{Chave}: Um campo ou valor associado a um registro, sobre o qual a ordenação é realizada. Por exemplo, ordenar uma lista de pessoas pela sua idade.
 \item \textbf{Ordenação interna vs. externa}:
    \begin{itemize}
        \item \textbf{Interna}: Ocorre quando todos os dados cabem na memória principal (RAM).
        \item \textbf{Externa}: Necessária quando os dados são muito volumosos e precisam ser mantidos em memória secundária (disco), sendo processados em partes.
    \end{itemize}
 \item \textbf{Estabilidade}: Um algoritmo é \textbf{estável} se ele preserva a ordem relativa original de registros com chaves iguais.
\end{itemize}
\end{frame}

\subsection{Eficiência}
\begin{frame}{Medindo a Eficiência}
\begin{itemize}
    \item A eficiência de um algoritmo de ordenação é tipicamente medida através de:
    \begin{itemize}
        \item \textbf{Complexidade de tempo}: Quantas operações o algoritmo executa em função do tamanho da entrada ($n$).
        \item \textbf{Complexidade de espaço}: Quanta memória adicional é necessária.
    \end{itemize}
    \item Para analisar o tempo, contamos as \textbf{operações críticas}:
    \begin{itemize}
        \item \textbf{Comparações}: O número de vezes que chaves são comparadas entre si.
        \item \textbf{Movimentações/Trocas}: O número de vezes que registros são movidos ou trocados de posição.
    \end{itemize}
    \item O \textbf{tamanho da entrada ($n$)} é o número de registros a serem ordenados.
\end{itemize}
\end{frame}

\section{Algoritmos Elementares}
\subsection{Ordenação por Flutuação (Bubble Sort)}
\begin{frame}{Bubble Sort: Estratégia}
\begin{itemize}
 \item Compara repetidamente pares de elementos adjacentes.
 \item Se estiverem na ordem errada, eles são trocados.
 \item A cada passagem completa, o maior elemento da porção não ordenada "flutua" para sua posição correta no final da lista.

\end{itemize}
\begin{center}
 \includegraphics[bb=0 0 1144 737,width=0.6\linewidth]{figs/EstrategiaBubbleSort.png}
 % EstrategiaBubbleSort.png: 1144x737 px, 72dpi, 40.36x26.00 cm, bb=0 0 1144 737
\end{center}
\end{frame}

\begin{frame}[fragile]{Bubble Sort: Algoritmo}
\begin{lstlisting}
funcao BubbleSort(A, tamanho)
  para i de 0 ate tamanho-2 faca
    // A cada passagem, o i-esimo maior elemento
    // eh colocado em sua posicao
    para j de 0 ate tamanho-i-2 faca
      // Compara elementos adjacentes
      se A[j] > A[j+1] entao
        troca(A[j], A[j+1])
      fim se
    fim para
  fim para
fim funcao
\end{lstlisting}
\begin{itemize}
    \item \textbf{Estabilidade}: Sim, a versão padrão do Bubble Sort é estável.

\end{itemize}
\end{frame}

\begin{frame}{Bubble Sort: Análise}
\begin{itemize}
    \item O número de comparações é fixo para a versão não otimizada.
    \item O laço externo roda $n-1$ vezes e o interno também, diminuindo a cada passo.
    \item \textbf{Comparações} (Todos os Casos):
    $$C(n) = (n-1) + (n-2) + \dots + 1 = \sum_{i=1}^{n-1} i = \frac{n(n-1)}{2}$$
    Portanto, a complexidade é sempre $O(n^2)$ para comparações.
    \item \textbf{Trocas}:
    \begin{itemize}
        \item \textbf{Melhor Caso} (lista ordenada): 0 trocas.
        \item \textbf{Pior Caso} (ordem reversa): $\frac{n(n-1)}{2}$ trocas.
    \end{itemize}
\end{itemize}

\end{frame}

\begin{frame}{Bubble Sort: Análise}
\begin{block}{Análise Assintótica}
\begin{itemize}
    \item Mesmo com otimizações, a complexidade geral do Bubble Sort é $O(n^2)$.
    \item É raramente usado em prática para problemas sérios, mas é um bom algoritmo para fins didáticos.
\end{itemize}
\end{block}
\end{frame}


\subsection{Ordenação por Inserção (Insertion Sort)}
\begin{frame}{Insertion Sort: Estratégia}
\begin{itemize}
 \item Percorre-se o vetor da esquerda para a direita. Cada elemento é "inserido" na sua posição correta dentro da sub-lista já ordenada à esquerda.
\end{itemize}
\begin{center}
 \includegraphics[bb=0 0 3508 2480,scale=0.05]{figs/EstrategioInsertSort.jpg}
 % EstrategioInsertSort.jpg: 3508x2480 px, 72dpi, 123.75x87.49 cm, bb=0 0 3508 2480
\end{center}
\end{frame}

\begin{frame}[fragile]
\begin{lstlisting}
funcao InsertionSort(A, tamanho)
  para j de 1 ate tamanho-1 faca
    chave = A[j]  // O elemento a ser inserido
    i = j - 1     // O fim da sub-lista ordenada

    // Move os elementos maiores que a chave
    enquanto (i >= 0) e (A[i] > chave) faca
      A[i+1] = A[i]
      i = i - 1
    fim enquanto

    A[i+1] = chave // Insere a chave
  fim para
fim funcao
\end{lstlisting}
\begin{itemize}
\item \textbf{Estabilidade}: Sim, o Insertion Sort é um algoritmo estável.
\end{itemize}
\end{frame}

\begin{frame}{Insertion Sort: Análise}
\begin{block}{Análise de Complexidade}
\begin{itemize}
    \item \textbf{Melhor Caso}: A lista já está ordenada.
    \begin{itemize}
        \item Comparações: O laço `enquanto` nunca executa seu corpo. $n-1$ comparações. Complexidade: $O(n)$.
        \item Movimentações: Nenhuma troca é necessária.
    \end{itemize}
    \item \textbf{Pior Caso}: A lista está em ordem reversa.
    \begin{itemize}
        \item Comparações: Para cada elemento $j$, são feitas $j-1$ comparações. Total: $\sum_{k=1}^{n-1} k = \frac{n(n-1)}{2}$. Complexidade: $O(n^2)$.
        \item Movimentações: Similar às comparações. Complexidade: $O(n^2)$.
    \end{itemize}
    \item \textbf{Caso Médio}: Complexidade de $O(n^2)$ para comparações e trocas.
\end{itemize}
\end{block}
\begin{itemize}
\item É muito eficiente para listas \textbf{pequenas} ou \textbf{"quase ordenadas"}.
\end{itemize}
\end{frame}


\subsection{Selection Sort (Ordenação por Seleção)}
\begin{frame}{Selection Sort: Estratégia}
 \begin{itemize}
  \item Divide a lista em duas partes: uma sub-lista ordenada à esquerda e uma sub-lista não ordenada à direita.
  \item Em cada iteração, encontra o \textbf{menor} elemento na sub-lista não ordenada.
  \item Troca esse menor elemento com o primeiro elemento da sub-lista não ordenada.
  \item A fronteira da sub-lista ordenada avança uma posição para a direita.
 \end{itemize}

\end{frame}

\begin{frame}{Estratégia}
 \begin{center}
 \includegraphics[bb=0 0 518 242,width=\textwidth]{figs/EstrategiaSelectionSort.jpg}
 % EstrategiaSelectionSort.jpg: 690x323 px, 96dpi, 18.26x8.55 cm, bb=0 0 518 242
\end{center}

\end{frame}

\begin{frame}[fragile]{Selection Sort: Algoritmo}
\begin{lstlisting}
funcao SelectionSort(A, tamanho)
  para i de 0 ate tamanho-2 faca
    // Encontra o indice do menor elemento restante
    min_idx = i
    para j de i+1 ate tamanho-1 faca
      se A[j] < A[min_idx] entao
        min_idx = j
      fim se
    fim para

    // Troca o menor encontrado com o primeiro
    // elemento da parte nao ordenada
    troca(A[i], A[min_idx])
  fim para
fim funcao
\end{lstlisting}
\begin{itemize}
    \item \textbf{Estabilidade}: Não, a implementação padrão não é estável.
\end{itemize}
\end{frame}

\begin{frame}
\begin{block}{Análise de Complexidade}
\begin{itemize}
    \item \textbf{Comparações}: O número de comparações é independente da ordem inicial dos dados.
    \begin{itemize}
        \item Os laços aninhados resultam em $\sum_{i=0}^{n-1} (n-1-i) = \frac{n(n-1)}{2}$ comparações.
        \item A complexidade é sempre $\Theta(n^2)$.
    \end{itemize}
    \item \textbf{Trocas}:
    \begin{itemize}
        \item O número de trocas é no máximo $n-1$, uma por iteração do laço externo.
        \item Isso faz do Selection Sort uma boa escolha quando o custo de mover os dados é muito alto.
    \end{itemize}
\end{itemize}
\end{block}
\begin{block}{Análise Assintótica}
 \begin{itemize}
    \item O tempo de execução é sempre $\Theta(n^2)$, independentemente da entrada.
 \end{itemize}
\end{block}
\end{frame}


\subsection{Ordenação Rápida (Quick Sort)}
\begin{frame}{Quick Sort: Características}
\begin{itemize}
 \item Um algoritmo recursivo baseado na estratégia de \textbf{dividir para conquistar}.
 \item Geralmente, é o mais rápido entre os algoritmos de ordenação baseados em comparação para dados em memória.
 \begin{itemize}
  \item \textbf{Caso Médio}: $\Theta(n \log n)$
  \item \textbf{Pior Caso}: $\Theta(n^2)$ (mas é raro com boas estratégias de pivô)
 \end{itemize}
 \item O núcleo do método é a rotina de \textbf{particionamento}.
  \begin{itemize}
  \item A partição rearranja a lista em torno de um elemento \textbf{pivô}, de modo que todos os elementos à esquerda do pivô são menores, e todos à direita são maiores.
 \end{itemize}
 \item \textbf{Conquistar}: As duas sub-listas (à esquerda e à direita do pivô) são ordenadas recursivamente.
 \item \textbf{Combinar}: Nenhum trabalho é necessário! A ordenação ocorre "no lugar" (in-place).
\end{itemize}
\end{frame}

\begin{frame}[fragile]
\begin{lstlisting}[mathescape=true]
funcao QuickSort(A, inicio, fim)
  se inicio < fim entao
    // P recebe o index do pivo que agora esta na posição correta
    p = Particao(A, inicio, fim)

    // Ordena recursivamente a sub-lista a esquerda de P
    QuickSort(A, inicio, p - 1)

    // Ordena recursivamente a sub-lista a direita de P
    QuickSort(A, p + 1, fim)
  fim se
fim funcao
\end{lstlisting}

\end{frame}

\begin{frame}[fragile]
\begin{lstlisting}[mathescape=true]
funcao Particao(A, inicio, fim)
  pivo = A[fim]//Estratégia: o ultimo elemento pivo
  // Começa no início do sub-array.
  i = inicio
  // 'j' percorre sub-array p/ encontrar elementos < P.
  para j de inicio ate fim - 1 faca
    // Se um elemento menor que o pivo e encontrado
    se A[j] < pivo entao
      troca(A[i], A[j])//Faz swap
      i = i + 1//avanca p/ proxima posicao.
    fim se
  fim para
  // Ao final,os elementos antes de 'i' sao menores que P.
  // Coloca P em sua posicao final correta, swap com A[i].
  troca(A[i], A[fim])
  // Retorna o index onde P foi colocado.
  retorne i
fim funcao
\end{lstlisting}
\begin{itemize}
    \item \textbf{Estabilidade}: Não, o Quick Sort geralmente não é estável.
\end{itemize}
\end{frame}

\begin{frame}[fragile]
\frametitle{Quicksort - Início da Partição}
\centering
\begin{tikzpicture}[
    node distance=8mm and 12mm,
    array/.style={
        matrix of nodes,
        nodes={draw, minimum size=8mm, anchor=center, font=\large},
        row sep=0pt,
        column sep=-\pgflinewidth
    },
    pivot/.style={fill=orange!40},
    less/.style={fill=green!30},
    greater/.style={fill=yellow!40},
    sorted/.style={fill=cyan!30},
    pointer/.style={-Stealth, blue, thick},
    legendbox/.style={draw, minimum size=6mm},
    info/.style={align=left, font=\small\itshape, text width=4cm}
]

% Legenda
\begin{scope}[start chain=legend_chain going right, node distance=2mm]
    \node[on chain, legendbox, less]{};
    \node[on chain, right=of legend_chain-1, anchor=west] {$<$pivô};
    \node[on chain, legendbox, greater, left=of legend_chain-2, xshift=2cm]{};
    \node[on chain, right=of legend_chain-3, anchor=west] {$>$pivô};
    \node[on chain, legendbox, pivot, left=of legend_chain-4, xshift=2cm]{};
    \node[on chain, right=of legend_chain-5, anchor=west] {Pivô};
    \node[on chain, legendbox, sorted, left=of legend_chain-6, xshift=2cm]{};
    \node[on chain, right=of legend_chain-7, anchor=west] {Posição final};
\end{scope}

% Estado inicial
\matrix[array, below=of legend_chain-1] (init) {
   5 & 3 & 8 & 4 & |[pivot]| 6 \\
};
\node[info, right=of init] (desc1) {Pivô = 6. Ponteiros i e j iniciam nas posições 0 e 1. \newline\newline Verifica A[j] $<$ Pivô, caso verdadeiro então faz swap com A[i]};
\node[below=3mm of init-1-1] (i0) {i};
\node[below=3mm of init-1-2] (j0) {j};
\draw[pointer] (i0.north) -- (init-1-1.south);
\draw[pointer] (j0.north) -- (init-1-2.south);

\end{tikzpicture}
\end{frame}

% --- Slide 2: Passo 1 ---
\begin{frame}[fragile]
\frametitle{Passo 1: j $<$ pivô}
\centering
\begin{tikzpicture}[array/.style={
        matrix of nodes,
        nodes={draw, minimum size=8mm, anchor=center, font=\large},
        row sep=0pt,
        column sep=-\pgflinewidth
    },
    pivot/.style={fill=orange!40},
    less/.style={fill=green!30},
    pointer/.style={-Stealth, blue, thick},
    swap/.style={-Stealth, red, thick, bend left=40},
    info/.style={align=left, font=\small\itshape, text width=4cm}
]

\matrix[array] (step1) {
   |[less]| 3 & 5 & 8 & 4 & |[pivot]| 6 \\
};
\node[info, right=of step1] (desc2) {Incrementa i e j.\newline\newline Verifica A[j] $<$ Pivô, caso verdadeiro então faz swap com A[i]};
\node[below=3mm of step1-1-2] (i1) {i};
\node[below=3mm of step1-1-3] (j1) {j};
\draw[pointer] (i1.north) -- (step1-1-2.south);
\draw[pointer] (j1.north) -- (step1-1-3.south);

\end{tikzpicture}
\end{frame}

\begin{frame}[fragile]
\frametitle{Passo 1: j $<$ pivô}
\centering
\begin{tikzpicture}[array/.style={
        matrix of nodes,
        nodes={draw, minimum size=8mm, anchor=center, font=\large},
        row sep=0pt,
        column sep=-\pgflinewidth
    },
    pivot/.style={fill=orange!40},
    less/.style={fill=green!30},
    pointer/.style={-Stealth, blue, thick},
    swap/.style={-Stealth, red, thick, bend left=40},
    info/.style={align=left, font=\small\itshape, text width=4cm}
]

\matrix[array] (step1) {
   |[less]| 3 & 5 & 8 & 4 & |[pivot]| 6 \\
};
\node[info, right=of step1] (desc2) {Incrementa i e j.\newline\newline Verifica A[j] $<$ Pivô, caso verdadeiro então faz swap com A[i] \newline\newline Incrementa j };
\node[below=3mm of step1-1-2] (i1) {i};
\node[below=3mm of step1-1-3] (j1) {j};
\draw[pointer] (i1.north) -- (step1-1-2.south);
\draw[pointer] (j1.north) -- (step1-1-3.south);

\end{tikzpicture}
\end{frame}

% --- Slide 3: Passo 2 ---
\begin{frame}[fragile]
\frametitle{Passo 2: j $>$ pivô}
\centering
\begin{tikzpicture}[array/.style={
        matrix of nodes,
        nodes={draw, minimum size=8mm, anchor=center, font=\large},
        row sep=0pt,
        column sep=-\pgflinewidth
    },
    pivot/.style={fill=orange!40},
    greater/.style={fill=yellow!40},
    pointer/.style={-Stealth, blue, thick},
    info/.style={align=left, font=\small\itshape, text width=4cm}
]

\matrix[array] (step2) {
   3 & 5 & |[greater]| 8 & 4 & |[pivot]| 6 \\
};
\node[info, right=of step2] (desc3) {Verifica A[j] $<$ Pivô, caso verdadeiro então faz swap com A[i]};
\node[below=3mm of step2-1-2] (i2) {i};
\node[below=3mm of step2-1-4] (j2) {j};
\draw[pointer] (i2.north) -- (step2-1-2.south);
\draw[pointer] (j2.north) -- (step2-1-4.south);

\end{tikzpicture}
\end{frame}

% --- Slide 4: Passo 3 ---
\begin{frame}[fragile]
\frametitle{Passo 3: j $<$ pivô}
\centering
\begin{tikzpicture}[array/.style={
        matrix of nodes,
        nodes={draw, minimum size=8mm, anchor=center, font=\large},
        row sep=0pt,
        column sep=-\pgflinewidth
    },
    pivot/.style={fill=orange!40},
    less/.style={fill=green!30},
    pointer/.style={-Stealth, blue, thick},
    swap/.style={-Stealth, red, thick, bend left=40},
    info/.style={align=left, font=\small\itshape, text width=4cm}
]

\matrix[array] (step3) {
   3 & |[less]| 4 & 8 & 5 & |[pivot]| 6 \\
};
\node[info, right=of step3] (desc4) {Incrementa i e j. \newline\newline Percorreu até final. Então swap (A[i],A[j]) $\rightarrow$ Troca o pivô com o elemento na posição i.)};
\node[below=3mm of step3-1-3] (i3) {i};
\node[below=3mm of step3-1-5] (j3) {j};
\draw[pointer] (i3.north) -- (step3-1-3.south);
\draw[pointer] (j3.north) -- (step3-1-5.south);

\end{tikzpicture}
\end{frame}

% --- Slide 5: Passo Final (Troca com pivô) ---
\begin{frame}[fragile]
\frametitle{Passo Final: Troca com o pivô}
\centering
\begin{tikzpicture}[array/.style={
        matrix of nodes,
        nodes={draw, minimum size=8mm, anchor=center, font=\large},
        row sep=0pt,
        column sep=-\pgflinewidth
    },
    sorted/.style={fill=cyan!30},
    info/.style={align=left, font=\small\itshape, text width=4cm},
    swap/.style={-Stealth, red, thick, bend left=40}
]

\matrix[array] (step4) {
   3 & 4 & |[sorted]| 6 & 5 & 8 \\
};
\node[info, right=of step4] (desc5) {Fim da partição. };

\end{tikzpicture}
\end{frame}

\begin{frame}{Quick Sort: Análise}
\begin{itemize}
 \item O desempenho do Quick Sort depende criticamente da \textbf{escolha do pivô}.
 \item O custo do particionamento é sempre linear, $\Theta(n)$.
\end{itemize}
\begin{description}
 \item[Pior Caso] Ocorre quando o pivô é consistentemente o menor ou o maior elemento. Isso cria partições desbalanceadas (uma com 0 elementos e outra com $n-1$). A complexidade se torna $O(n^2)$.
 \item[Melhor Caso] O pivô é sempre o elemento mediano, dividindo a lista em duas metades perfeitas. A complexidade é $O(n \log n)$.
 \item[Caso Médio] Para entradas aleatórias, o desempenho é $O(n \log n)$.
\end{description}

\end{frame}


\begin{frame}{Quick Sort: Análise}

\begin{block}{Como evitar o pior caso?}
 \begin{itemize}
    \item \textbf{Pivô Aleatório}: Escolher um elemento aleatório como pivô torna o pior caso extremamente improvável.
    \item \textbf{Mediana de Três}: Escolher o pivô como a mediana do primeiro, meio e último elemento da lista.
 \end{itemize}
\end{block}
\end{frame}

\subsection{Merge Sort (Ordenação por Intercalação)}
\begin{frame}{Merge Sort: Estratégia}
\begin{itemize}
 \item Também utiliza a abordagem de \textbf{dividir para conquistar}.
 \item O algoritmo segue três passos:
 \begin{description}
  \item[Dividir] O vetor é dividido recursivamente ao meio até que cada sub-vetor contenha apenas um elemento (que, por definição, está ordenado).
  \item[Conquistar] Esta fase é implícita; os vetores de um elemento já estão "conquistados" (ordenados).
  \item[Combinar] Os sub-vetores ordenados são \textbf{intercalados} (merged) de volta, de dois em dois, para produzir novos sub-vetores ordenados maiores, até que a lista inteira esteja combinada e ordenada.
 \end{description}
\end{itemize}
\end{frame}

\begin{frame}{}
\begin{center}
\begin{tikzpicture}[
    level 1/.style={sibling distance=5cm, level distance=1.0cm},
    level 2/.style={sibling distance=2.0cm},
    level 3/.style={sibling distance=1.0cm},
    level 4/.style={sibling distance=1cm},
    down/.style={->, >=Latex},
    up/.style={<-, >=Latex},
    arrnode/.style={font=\ttfamily}
]
\node[arrnode] {77 44 99 66 33 55 88 22}
    child { node[arrnode] {77 44 99 66}
        child { node[arrnode] {77 44}
            child { node[arrnode] {77} }
            child { node[arrnode] {44} }
        }
        child { node[arrnode] {99 66}
            child { node[arrnode] {99} }
            child { node[arrnode] {66} }
        }
    }
    child { node[arrnode] {33 55 88 22}
        child { node[arrnode] {33 55}
            child { node[arrnode] {33} }
            child { node[arrnode] {55} }
        }
        child { node[arrnode] {88 22}
            child { node[arrnode] {88} }
            child { node[arrnode] {22} }
        }
    };

    % Merge arrows and results
    \uncover<2->{
    \node[arrnode, sorted] at (-3.75, -4.25) {44 77};
    \node[arrnode, sorted] at (-1.25, -4.25) {66 99};
    \node[arrnode, sorted] at (1.25, -4.25) {33 55};
    \node[arrnode, sorted] at (3.75, -4.25) {22 88};
    \draw[up, thick, blue] (-3.75, -4) -- (-3.75, -3.5);
    \draw[up, thick, blue] (-1.25, -4) -- (-1.25, -3.5);
    \draw[up, thick, blue] (1.25, -4) -- (1.25, -3.5);
    \draw[up, thick, blue] (3.75, -4) -- (3.75, -3.5);
    }

    \uncover<3->{
    \node[arrnode, sorted] at (-2.5, -5) {44 66 77 99};
    \node[arrnode, sorted] at (2.5, -5) {22 33 55 88};
    \draw[up, thick, blue] (-2.5, -4.75) -- (-2.5, -4.25);
    \draw[up, thick, blue] (2.5, -4.75) -- (2.5, -4.25);
    }
    \uncover<4->{
    \node[arrnode, sorted] at (0, -5.75) {22 33 44 55 66 77 88 99};
    \draw[up, thick, blue] (0, -5.5) -- (0, -5);
    }

\end{tikzpicture}
\end{center}
\end{frame}

\begin{frame}[fragile]{Merge Sort: Algoritmo}
\begin{lstlisting}
procedimento MergeSort(A, esq, dir)
  se esq < dir entao
    meio = floor((esq + dir) / 2)
    MergeSort(A, esq, meio)
    MergeSort(A, meio + 1, dir)
    Intercalar(A, esq, meio, dir)
  fim se
fim procedimento
\end{lstlisting}
\end{frame}


\begin{frame}[fragile]{}
\begin{lstlisting}
procedimento Intercalar(A, esq, meio, dir)
  // Criar vetores temporarios
  n1 = meio - esq + 1
  n2 = dir - meio
  E[0..n1-1], D[0..n2-1]
  // Copiar dados para os vetores temporarios
  para i de 0 ate n1-1 faca E[i] = A[esq + i]
  para j de 0 ate n2-1 faca D[j] = A[meio + 1 + j]
  // Intercalar de volta para A
  i=0, j=0, k=esq
  enquanto i < n1 e j < n2 faca
    se E[i] <= D[j] entao A[k++] = E[i++]
    senao A[k++] = D[j++]
  fim enquanto
  // Copiar o restante, se houver
  enquanto i < n1 faca A[k++] = E[i++]
  enquanto j < n2 faca A[k++] = D[j++]
fim procedimento
\end{lstlisting}
\end{frame}

\begin{frame}{Merge Sort: Análise}
\begin{block}{Análise de Complexidade}
\begin{itemize}
 \item A divisão do vetor em metades cria uma árvore de recursão de altura $\log_2 n$.
 \item Em cada nível da árvore, a operação de intercalar processa todos os $n$ elementos.
 \item Portanto, o tempo de execução é o produto da altura da árvore pelo trabalho em cada nível.
 \item \textbf{Tempo de Execução}: $\Theta(n \log n)$ para \textbf{todos os casos} (melhor, médio e pior). O comportamento do algoritmo é consistente e não depende da ordem inicial dos dados.
\end{itemize}
\end{block}

\end{frame}


\begin{frame}{Merge Sort: Análise}

\begin{block}{Análise de Espaço}
\begin{itemize}
    \item O ponto fraco do Merge Sort é sua necessidade de memória auxiliar.
    \item A operação de intercalação requer um vetor temporário de tamanho $n$.
    \item \textbf{Complexidade de Espaço}: $O(n)$.
    \item \textbf{Estabilidade}: Sim, o Merge Sort é um algoritmo estável.
\end{itemize}
\end{block}
\end{frame}

\subsection{Trabalho 1 - Ordenação}
\begin{frame}{Trabalho Prático}
\begin{block}{Exercício 1: Comparando Algoritmos Elementares}
\begin{itemize}
 \item Implemente o \textbf{Insertion Sort} e o \textbf{Bubble Sort}.
 \item Para cada um, meça o \textbf{tempo de execução} e o \textbf{número de movimentações/trocas}.
 \item Execute os experimentos com 3 conjuntos de dados gerados aleatoriamente, com tamanhos:
 \begin{itemize}
 \item $n = 10.000$
 \item $n = 50.000$
 \item $n = 100.000$
 \end{itemize}
 \item Analise e compare os resultados.
\end{itemize}
\end{block}


\end{frame}

\begin{frame}{Trabalho Prático}


\begin{block}{Exercício 2: Otimizando o Quick Sort}
\begin{itemize}
 \item Usando como base o algoritmo Quick Sort apresentado, implemente duas estratégias de escolha de pivô: \textbf{"mediana-de-três"} e \textbf{"versão aleatorizada"}.
 \item Compare o desempenho (tempo de execução) das novas versões com a implementação original (pivô fixo) para um conjunto de dados aleatórios grande (e.g., $n = 1.000.000$).
 \item Discuta se as estratégias melhoraram o desempenho e por quê.
\end{itemize}
\end{block}
\end{frame}

\section{Referências Bibliográficas}
\begin{frame}[t,allowframebreaks]
\frametitle{Referências}
\begin{thebibliography}{9}
    \bibitem{A03} Cormen, T. H., Leiserson, C. E., Rivest, R. L., \& Stein, C. (2009). \textit{Introduction to Algorithms (3rd ed.)}. MIT Press.
    \bibitem{A04} Ascencio, A. F. G., \& Araujo, G. S. (2010). \textit{Estrutura de dados: Algoritmos, análise da complexidade e implementações em Java e C++}. Pearson Prentice Hall.
    \bibitem{A06} Szwarcfiter, J. L., \& Markenzon, L. (2015). \textit{Estruturas de dados e seus algoritmos (3a ed.)}. LTC.
    \bibitem{A01} Lee, K. D., \& Hubbard, S. (2015). Data Structures and Algorithms with Python. Springer.
    \bibitem{A02} Hubbard, J. (2007). \textit{Schaum’s Outline of Data Structures with Java}. McGraw-Hill.
\end{thebibliography}
\end{frame}

\section{}
\setbeamertemplate{logo}{}
\begin{frame}
\titlepage
\end{frame}

\end{document}
